% 唐虹 -- 中文简历
% Source maintained for the homepage CV page.

\documentclass[11pt,a4paper]{ctexart}
\usepackage[margin=1.8cm]{geometry}
\usepackage[hidelinks]{hyperref}
\usepackage{enumitem}
\setlist[itemize]{leftmargin=1.2em,itemsep=2pt,topsep=2pt}
\pagenumbering{gobble}

\begin{document}

\begin{center}
{\Large\textbf{唐虹}}\\
助理教授（研究），建筑环境与能源工程学系\\
香港理工大学\\
\href{mailto:honghong.tang@polyu.edu.hk}{honghong.tang@polyu.edu.hk} \quad
\href{https://scholar.google.com.hk/citations?user=nXW7KrwAAAAJ&hl=en}{Google Scholar} \quad
\href{https://www.researchgate.net/profile/Hong-Tang-17}{ResearchGate}
\end{center}

\section*{研究方向}
电网友好建筑；建筑能源柔性与需求响应；人工智能及机器学习在能源系统中的应用；算电协同及数据中心精密空调优化控制；综合能源系统与虚拟电厂。

\section*{教育背景}
\begin{itemize}
\item 香港理工大学，建筑环境与能源工程，博士，2018--2023。
\item 清华大学，建筑环境与能源应用工程，学士，2014--2018。
\end{itemize}

\section*{工作经历}
\begin{itemize}
\item 助理教授（研究），香港理工大学，2025年至今。
\item 博士后研究员，香港理工大学，2023--2025。
\item 研究助理，香港理工大学，2022--2023。
\end{itemize}

\section*{代表性论文}
\begin{itemize}
\item Hong Tang 等，Multi-Timescale Rolling Optimization Coordinating Heterogeneous EV and Building Load Flexibility in Community Energy Systems with NGBoost-Based Renewable Forecast，\emph{Applied Energy}，已接收，2026。
\item Hong Tang 等，Risk-aware optimal dispatch of resource aggregators integrating NGBoost-based probabilistic renewable prediction and bi-level building flexibility engagement，\emph{Engineering}，53:76--89，2025。
\item Hong Tang 等，Unlocking Building Energy Flexibility: Multi-Layered Pricing Framework Design from Distribution-Level to P2P Markets，\emph{The Innovation Energy Use}，1(2):100019，2025。
\item Hong Tang and Shengwei Wang，Game-theoretic optimization of demand-side flexibility engagement considering the perspectives of different stakeholders and multiple flexibility services，\emph{Applied Energy}，332:120550，2023。
\item Hong Tang, Shengwei Wang and Hangxin Li，Flexibility categorization, sources, capabilities and technologies for energy-flexible and grid-responsive buildings，\emph{Energy}，219:119598，2021。
\end{itemize}

\section*{服务与荣誉}
\begin{itemize}
\item \emph{The Innovation Energy Use} 青年编委。
\item 担任 \emph{Energy}、\emph{Building Simulation}、\emph{Buildings}、\emph{Architecture}、\emph{Journal of Physics: Conference Series}、\emph{Journal of Building Design and Environment} 客座编辑。
\item 最佳汇报奖，The 11th Asia Conference on Power and Electrical Engineering，2026。
\item 最佳墙报奖，The 4th Asia Conference of International Building Performance Simulation Association，2018。
\end{itemize}

\vfill
\noindent 最近更新：2026年5月27日

\end{document}
